\section{Marco teórico}
\subsection{Diseño de mezclas y ley de Abrams}
El diseño de mezclas de concreto se apoya históricamente en relaciones empíricas entre la resistencia a compresión y la relación agua/cemento (a/c). La ley de Abrams propone una relación inversa entre ambas magnitudes, que en forma log-lineal puede expresarse como:
\[
f'_c = K \cdot \left(\frac{w}{c}\right)^{-n}
\quad\text{o}\quad
\log f'_c = \log K - n \cdot \log\!\left(\frac{w}{c}\right),
\]
donde $f'_c$ es la resistencia característica, $K$ y $n$ son parámetros dependientes de materiales y condiciones de curado. Métodos normativos (p. ej., ACI 211.1) emplean esta y otras relaciones para estimar dosificaciones iniciales que luego se ajustan con ensayos. En DIAPCE, la relación a/c es un nodo clave del flujo de decisión y se restringe a valores presentes en el conjunto experimental, reduciendo la incertidumbre del usuario al seleccionar solo combinaciones observadas.

\subsection{Efecto de condiciones ambientales (temperatura y humedad)}
La cinética de hidratación del cemento y, por tanto, el desarrollo de la resistencia, dependen de la temperatura y el estado higrotérmico. Temperaturas más altas aceleran la hidratación inicial pero pueden reducir resistencias finales si hay pérdida de humedad; ambientes más secos favorecen gradientes de humedad y retracción. Por esto, DIAPCE incorpora explícitamente:
\begin{itemize}
    \item \textbf{Temperatura ambiente (°C)}: primer filtro del flujo; acota el subconjunto de registros.
    \item \textbf{Humedad relativa (\%)}: segundo filtro; condiciona las opciones de relación a/c.
\end{itemize}
Esta cascada emula un \textit{árbol de decisión} que condiciona las alternativas posteriores a partir de parámetros ambientales, manteniendo la coherencia con los datos observados.

\subsection{Aditivos químicos y códigos P0–PP3}
Los aditivos modifican trabajabilidad, tiempo de fraguado y desarrollo de resistencia (p. ej., reductores de agua, superplastificantes, acelerantes, retardantes; ASTM C494). En el dataset se normalizan bajo códigos \texttt{P0}–\texttt{PP3}, lo que permite:
\begin{itemize}
    \item \textbf{Selección guiada}: se muestran solo aditivos válidos para la combinación temperatura–humedad–a/c.
    \item \textbf{Trazabilidad}: el código simplifica el JOIN con tablas técnicas (\texttt{tipos\_aditivo}, \texttt{productos}, \texttt{aditivos}).
\end{itemize}

\subsection{Evolución temporal de la resistencia (7, 14 y 28 días)}
La resistencia se reporta típicamente a 7, 14 y 28 días (ASTM C39). A 28 días se toma como referencia de diseño. DIAPCE estima las resistencias $\{f_7, f_{14}, f_{28}\}$ como promedios empíricos del subconjunto filtrado por condiciones ambientales, relación a/c y aditivo:
\[
\hat f_{t} \;=\; \operatorname{avg}\{\, f_{t,i}\,:\, (T_i,H_i,\mathrm{a/c}_i,\mathrm{aditivo}_i)\in \text{selección}\,\},\quad t\in\{7,14,28\}.
\]
Este enfoque \textit{data-driven} privilegia la reproducibilidad y la coherencia con el comportamiento observado en laboratorio del conjunto local (1874 registros).

\subsection{Modelado empírico y enfoques de aprendizaje}
La literatura reciente explora desde regresiones multivariadas hasta redes neuronales para predecir $f'_c$ incorporando variables como a/c, tipo de cemento, granulometría, aditivos y condiciones de curado. DIAPCE adopta en esta versión un \textbf{modelo empírico de agregación} (promedios condicionados) por su:
\begin{itemize}
    \item \textbf{Transparencia}: resultados explicables y trazables a registros concretos.
    \item \textbf{Robustez} ante datos ruidosos cuando se hace estratificación adecuada.
    \item \textbf{Compatibilidad} con decisiones en cascada y validaciones de rango (resistencia objetivo 24–57 MPa con decimales).
\end{itemize}
La arquitectura y el esquema de datos dejan abierta la integración futura de modelos ML (p. ej., regresión regularizada o \textit{tree ensembles}) entrenados sobre las mismas vistas SQL.

\subsection{Sistemas de soporte a decisiones y consultas relacionales}
DIAPCE actúa como un \textit{asistente de decisión} local. Su núcleo es una base de datos SQLite con:
\begin{itemize}
    \item \textbf{Normalización y llaves foráneas} para garantizar integridad (proyectos, mezclas, materiales, aditivos, resultados de concreto).
    \item \textbf{Índices} (p. ej., \texttt{idx\_busqueda\_resistencia}) para acelerar las consultas condicionales.
    \item \textbf{Vistas/consultas} que implementan el árbol de decisión: 
    temperatura $\rightarrow$ humedad $\rightarrow$ relación a/c $\rightarrow$ aditivo, 
    y las funciones de predicción $\{\hat f_7,\hat f_{14},\hat f_{28}\}$.
\end{itemize}
Este patrón \textit{relational-first} asegura reproducibilidad de resultados y portabilidad sin dependencia de red.

\subsection{Arquitectura móvil y visualización}
La aplicación se implementa en Flutter (Dart), con almacenamiento local (\texttt{sqflite}) y visualizaciones con \texttt{fl\_chart}:
\begin{itemize}
    \item \textbf{UI guiada}: formularios con validadores que restringen a intervalos y opciones presentes en el dataset.
    \item \textbf{Gráficos}: pastel para composición de mezcla y línea para la evolución de resistencia, apoyando interpretación técnica.
    \item \textbf{Proyectos por usuario}: aislamiento lógico que facilita trazabilidad y organización de diseños.
\end{itemize}

\subsection{Síntesis conceptual aplicada en DIAPCE}
Integrando los elementos anteriores, el marco teórico de DIAPCE se puede resumir en tres pilares:
\begin{enumerate}
    \item \textbf{Fundamento de materiales}: relación a/c (ley de Abrams), influencia higrotérmica, y efecto de aditivos en la hidratación y trabajabilidad.
    \item \textbf{Evidencia experimental local}: predicciones y opciones guiadas extraídas de un dataset real de 1874 registros, con estratificación por condiciones.
    \item \textbf{Soporte computacional}: consultas relacionales optimizadas, decisión en cascada y visualizaciones que transforman datos en evidencia accionable.
\end{enumerate}

\subsection{Trabajos relacionados y estándares de referencia}
Como línea base se consideran los métodos ACI para diseño de mezclas (ACI 211.1) y normas de ensayo (ASTM C39 para compresión, ASTM C494 para aditivos), junto con literatura de optimización y predicción de resistencia (p. ej., relaciones empíricas y regresión multivariable). DIAPCE se posiciona como un \textit{asistente data-driven} local, complementario a softwares comerciales, con énfasis en adaptabilidad regional y trazabilidad de decisiones.